% ============================================================
%  Capitolo 3 — Architettura del sistema CDSS neuro-simbolico
%  Tesi Triennale Francesco Zamar — UniTS 2026
%  Versione finale (rifatta da cap3 v4 + estratti architetturali cap4)
% ============================================================
\chapter{Architettura del sistema}
\label{ch:progettazione}

Il \Cref{ch:background} ha definito tre strumenti: un linguaggio per
rappresentare l'incertezza, una pipeline per ancorare le risposte ai
documenti, un confine concettuale fra ragionamento simbolico e generativo.
La domanda di progetto è come comporli in un sistema unico che rimanga
clinicamente sicuro, cioè in cui i fallimenti possibili del componente
generativo, che è intrinsecamente non verificabile per costruzione
\cite{ji2023hallucination, thirunavukarasu2023llm}, non possano mai
ribaltare una decisione di rimborsabilità. La risposta che questo capitolo
sviluppa è la \emph{separazione neuro-simbolica delle responsabilità}: la
correttezza normativa è interamente affidata a un rule engine deterministico
ispezionabile, mentre il modello generativo riceve la decisione già calcolata
e ha l'unico compito di articolarne la motivazione in linguaggio clinico
ancorato alle Note. La discussione procede dai requisiti che il sistema deve
soddisfare, attraverso il principio architetturale che li governa, fino alla
descrizione analitica di ciascun componente e del contratto che li lega.

% ─────────────────────────────────────────────────────────────
\section{Requisiti e separazione neuro-simbolica}
\label{sec:requisiti}

Il sistema viene progettato per operare su quattro Note AIFA (Agenzia Italiana del Farmaco): la \notezero\
(inibitori di pompa protonica), la \notetredici\ (ipolipemizzanti), la \notesessanta\
(antinfiammatori non steroidei) e la \notenova\ (anticoagulanti orali nella
fibrillazione atriale non valvolare).
Il corpus di valutazione comprende \goldn{} casi di riferimento raccolti
per coprire tutti i percorsi decisionali ammissibili.

Sul piano funzionale, il sistema deve soddisfare sette obiettivi distinti,
ciascuno con un criterio di accettazione misurabile.
Il primo è la verifica formale della rimborsabilità: il sistema deve produrre
un esito appartenente all'insieme
\[
  \mathcal{D} = \{\textsc{rimborsabile},\;
  \textsc{non\_rimborsabile},\;
  \textsc{non\_determinabile},\;
  \textsc{routed}\}
\]
con Macro F1 pari a~1,0 sull'intero corpus di riferimento.
Il secondo è la spiegazione in linguaggio naturale: ogni risposta deve contenere
almeno un'evidenza che citi esplicitamente il PDF sorgente, il numero di pagina
e un estratto verbatim dalla norma applicata.
Il terzo è la citazione di estratti normativi verificabili: la copertura
delle citazioni deve essere del 100\,\%, nel senso che ogni decisione deve
associare almeno un'ancora \emph{file\,+\,page\,+\,excerpt} al testo PDF.
Il quarto requisito riguarda la gestione dell'incertezza: l'assenza di un
campo clinico nel profilo paziente si deve propagare come
\texttt{UNKNOWN} nella logica a tre valori, senza imporre valori di default
né degradare silenziosamente verso \texttt{FALSE}, il che trasformerebbe
un caso indeterminato in un rifiuto non fondato.
Il quinto requisito è la modalità what-if: deve essere possibile modificare
un campo del profilo paziente e rieseguire la sola valutazione simbolica
senza reinvocare il layer linguistico, qualora la decisione non sia cambiata.
Il sesto è la segnalazione strutturata dei campi mancanti: in caso di esito
\decision{NON\_DETERMINABILE}, la risposta deve riportare la lista esatta
delle variabili cliniche la cui assenza ha causato l'indeterminazione,
in modo che il clinico sappia quali informazioni raccogliere per risolvere
l'ambiguità.
Il settimo e ultimo requisito funzionale è la tracciabilità e l'audit
della catena decisionale: l'oggetto di risposta deve includere la sequenza
ordinata delle regole valutate, con l'identificatore, il valore Kleene
e l'ancora normativa per ciascuna, esportabile ai fini di verifica regolatoria.

Sul piano non funzionale, quattro proprietà condizionano l'intera architettura.
Il \emph{determinismo} delle decisioni cliniche richiede che,
a parità di input, l'esito simbolico sia sempre identico, indipendentemente
dall'ordine di valutazione o dallo stato interno dei processi;
questo esclude qualsiasi fonte di stocasticità nel layer decisionale.
La \emph{verificabilità per-caso} impone che ogni singola decisione possa
essere rieseguita e auditata in isolamento, senza bisogno di avviare
l'intero sistema: il rule engine simbolico deve essere esercitabile
autonomamente sull'insieme di riferimento anche quando il layer RAG (\emph{Retrieval-Augmented Generation}) non è attivo.
La \emph{riproducibilità canonica} della valutazione richiede che
le metriche siano misurabili su un ambiente riproducibile da zero:
wipe del database vettoriale, re-ingestione da PDF puliti e valutazione
end-to-end, con dipendenze Python fissate tramite hash crittografici.
Infine, la \emph{latenza accettabile per uso clinico interattivo} non
impone requisiti sub-secondo (l'interazione con un medico prescrivente
tollera qualche secondo di elaborazione), ma esclude tempi di risposta
dell'ordine dei minuti che renderebbero il sistema inutilizzabile nella
pratica quotidiana.
Un vincolo trasversale di progetto è il deployment interamente locale:
nessun dato clinico deve lasciare la macchina su cui il sistema è installato,
in coerenza con i requisiti di riservatezza dei dati sanitari e con
le limitazioni infrastrutturali degli ambienti ospedalieri italiani.

La separazione fra decisione e spiegazione non è soltanto una scelta
architetturale: è un vincolo di sicurezza.
Un large language model, per quanto accurato nelle sue previsioni medie,
non offre garanzie formali sulla coerenza fra il suo output e una specifica
normativa scritta~\cite{ji2023hallucination}.
L'allucinazione, la produzione di affermazioni plausibili ma non fondate,
è una proprietà intrinseca dei modelli autoregressivi, non un difetto
che la selezione del modello o l'aumento della scala risolvono
completamente~\cite{thirunavukarasu2023llm}.
Un rule engine, al contrario, è verificabile per costruzione:
ogni decisione è la conseguenza deterministica di una catena di regole
esplicite e ispezionabili.
Per questo motivo, nel sistema proposto, l'unica responsabilità del
modello generativo è la produzione del testo esplicativo, mentre
la determinazione della rimborsabilità è interamente affidata
al layer simbolico.

Il trade-off è reale e va riconosciuto onestamente.
Il layer simbolico richiede che le regole normative siano trascritte
manualmente in un formato strutturato, un lavoro che introduce la possibilità
di errori di codifica non presenti nel testo PDF originale;
il sistema di test e i casi di riferimento mitigano questo rischio
ma non lo eliminano del tutto.
Il layer linguistico introduce latenza aggiuntiva, stocasticità residua
anche con temperatura fissata a zero, e la necessità di un controllo
post-generazione che verifichi la coerenza fra il testo prodotto e la
decisione deterministica già calcolata.
Entrambi i layer, inoltre, hanno footprint computazionale e devono
coesistere su hardware commodity.
La scelta di separarli architetturalmente trasforma questi rischi in
dipendenze gestibili: il rule engine può essere testato esaustivamente
in isolamento, e il layer linguistico può essere sostituito o aggiornato
senza toccare la logica decisionale.

\begin{definition}[Architettura neuro-simbolica con separazione delle responsabilità]
\label{def:separazione}
Il sistema $\mathcal{S}$ è detto \emph{neuro-simbolico con separazione
delle responsabilità} se e solo se esiste una partizione
$\mathcal{S} = \mathcal{S}_{\text{sym}} \cup \mathcal{S}_{\text{neu}}$
tale che: (i)~$\mathcal{S}_{\text{sym}}$ è verificabile formalmente
su un test set $\mathcal{T}$ senza invocare alcun componente di
$\mathcal{S}_{\text{neu}}$; (ii)~$\mathcal{S}_{\text{neu}}$ riceve
come input l'output già calcolato di $\mathcal{S}_{\text{sym}}$
e non può modificarlo; (iii)~il contratto fra i due strati è espresso
da un tipo strutturato con schema validato a runtime.
\end{definition}

Il secondo principio architetturale che governa il design è il
\emph{fail-fast}: ogni fase di validazione deve fallire immediatamente
e in modo esplicito al primo invariante violato.
Nel rule engine questo si traduce in fasi ordinate con uscite di
short-circuit esplicite; nel layer RAG si traduce nel controllo
post-generazione che rileva le contraddizioni fra la decisione
deterministica e il testo prodotto dall'LLM prima che raggiungano
il clinico.
Uno stato corrotto silenzioso è più pericoloso di un errore esplicito
in un sistema che supporta decisioni cliniche.

I principi di \emph{single source of truth} e di \emph{traceability}
completano il quadro: le regole normative hanno un'unica rappresentazione
(i file YAML versionati), e ogni decisione deve essere ricondotta a una
catena esplicita e riproducibile di passi, ognuno con l'ancora nel PDF
che ne è la fonte autorizzativa.

% ─────────────────────────────────────────────────────────────
\section{Architettura complessiva}
\label{sec:architettura}

Il sistema è articolato in sei componenti principali, organizzati in
tre livelli funzionali: il livello decisionale simbolico, il livello
esplicativo neurale e il livello di presentazione.
La \Cref{fig:architettura} ne illustra la \emph{deployment view},
mostrando i confini di processo, i protocolli di comunicazione e
il flusso di dati sia nella fase online di valutazione sia nella
fase offline di ingestione dei PDF.

\begin{figure}[ht]
\centering
\resizebox{\textwidth}{!}{%
\begin{tikzpicture}[
  box/.style={
    rectangle, draw=black!70, thick,
    minimum width=3.2cm, minimum height=1.1cm,
    align=center, font=\small
  },
  store/.style={
    rectangle, draw=black!50, dashed, thick,
    minimum width=3.2cm, minimum height=0.9cm,
    align=center, font=\small
  },
  arrow/.style={-Stealth, thick},
  dasharrow/.style={-Stealth, thick, dashed},
  lbl/.style={font=\scriptsize, fill=white, inner sep=1pt},
  node distance=1.8cm and 3.2cm
]

%% Riga superiore — Demo
\node[box] (demo)
  {Demo Streamlit\\{\footnotesize porta 8501}};

%% Riga centrale — Orchestratore
\node[box, below=of demo] (orch)
  {\texttt{CDSSOrchestrator}\\{\footnotesize porta 8001 · FastAPI}};

%% Colonna sinistra — Rule Engine
\node[box, left=of orch] (re)
  {Rule Engine\\{\footnotesize porta 8000 · FastAPI}};

%% Storage collegato al Rule Engine
\node[store, above=1.2cm of re] (yaml)
  {Regole YAML\\{\footnotesize + catalogo principi attivi}};
\node[store, below=1.2cm of re] (pdf)
  {PDF Store\\{\footnotesize 4 Note AIFA}};

%% Colonna destra — Ollama
\node[box, right=of orch] (ollama)
  {Ollama\\{\footnotesize porta 11434}\\{\footnotesize \llmmodel}};

%% ChromaDB — sotto l'orchestratore
\node[store, below=1.4cm of orch] (chroma)
  {ChromaDB\\{\footnotesize embedded · 4 collezioni}};

%% Frecce online
\draw[arrow] (demo.south) -- node[lbl,right]{\texttt{CDSSRequest}} (orch.north);
\draw[arrow] (orch.north) -- node[lbl,left]{\texttt{CDSSResponse}} (demo.south);

\draw[arrow] (orch.west) -- node[lbl,above]{\texttt{/evaluate}} (re.east);
\draw[arrow] (re.east)   -- node[lbl,below]{\texttt{EvaluationResult}} (orch.west);

\draw[arrow] (yaml.south) -- (re.north);

\draw[arrow] (orch.east) -- node[lbl,above]{\textit{prompt}} (ollama.west);
\draw[arrow] (ollama.west) -- node[lbl,below]{\textit{spiegazione}} (orch.east);

\draw[arrow] (orch.south) -- node[lbl,right]{\texttt{retrieve()}} (chroma.north);
\draw[arrow] (chroma.north) -- node[lbl,left]{\textit{chunks}} (orch.south);

%% Freccia offline (ingestione)
%% pos=0.35: a metà curva l'etichetta copriva il testo del nodo ChromaDB.
\draw[dasharrow] (pdf.east) to[out=0, in=220]
  node[lbl, below, pos=0.35]{\texttt{ingest\_v2} (offline)} (chroma.west);

\end{tikzpicture}%
}
\caption{Deployment view dell'architettura a componenti. Le frecce continue
  rappresentano il flusso online durante la valutazione (REST/HTTP tra processi
  separati; filesystem per ChromaDB e PDF); la freccia tratteggiata indica
  la fase offline di ingestione: i PDF normativi vengono estratti, segmentati
  e indicizzati nelle quattro collezioni ChromaDB~\cite{chromadb}.}
\label{fig:architettura}
\end{figure}

Il Rule Engine risiede in un microservizio autonomo esposto sulla porta~8000
tramite FastAPI~\cite{fastapi} e Pydantic~\cite{pydantic}.
La scelta di isolarlo come processo separato risponde direttamente al requisito
di verificabilità per-caso: l'intero insieme di riferimento di \goldn{} casi
può essere esercitato senza che Ollama sia in esecuzione, riducendo
drasticamente il tempo di ogni ciclo test-fix-verify.
Il costo di questa scelta è un'interfaccia REST intermedia che introduce
una latenza aggiuntiva nell'ordine dei millisecondi, trade-off accettabile
dato che l'interazione clinica non impone requisiti sub-secondo.

L'Orchestratore, in esecuzione sulla porta~8001, costituisce il componente
di integrazione: interroga il Rule Engine tramite HTTP/REST~\cite{fielding2000rest},
raccoglie i chunk pertinenti da ChromaDB e compone il prompt strutturato per
Ollama~\cite{ollama2024}.
Il canale REST fra Rule Engine e Orchestratore è formalmente tipizzato tramite
Pydantic: ogni scambio di dati attraversa uno schema validato, eliminando
la categoria di errori dovuta a contratti impliciti.

ChromaDB opera in modalità embedded con persistenza su filesystem locale,
eliminando i rischi di disponibilità connessi al teorema CAP (\emph{Consistency, Availability, Partition tolerance})~\cite{brewer2000cap}
che affliggono le soluzioni distribuite.
Il deployment è interamente locale: nessun dato clinico lascia la macchina
su cui il sistema è installato, in coerenza con i requisiti di riservatezza
dei dati sanitari e con le limitazioni infrastrutturali degli ambienti
ospedalieri italiani.
La containerizzazione Docker, coordinata da un \texttt{docker-compose.yml}
che fissa le versioni di tutti i servizi, garantisce la riproducibilità
dell'ambiente di esecuzione indipendentemente dall'hardware host.

La demo interattiva Streamlit (porta~8501) fornisce un'interfaccia grafica
per la valutazione manuale dei casi: riceve il profilo paziente, invoca
l'Orchestratore e presenta la decisione simbolica insieme alla spiegazione
testuale con le evidenze citate.

% ─────────────────────────────────────────────────────────────
\section{Il layer simbolico: rule engine}
\label{sec:layer-simbolico}

Il Rule Engine implementa la componente simbolica del sistema.
La sua responsabilità esclusiva è la determinazione dell'esito di
rimborsabilità: dato un profilo paziente e una prescrizione, restituisce
un oggetto \texttt{EvaluationResult} strutturato che incapsula la decisione,
la traccia di copertura per ogni regola valutata, i flag clinici derivati
e la lista dei campi mancanti che hanno causato eventuali indeterminazioni.

\subsection{Modello dei dati e schema YAML}
\label{subsec:modello-dati}

Il profilo paziente è rappresentato dal tipo \texttt{PatientProfile},
che aggrega tre sotto-profili: dati anagrafici (età, sesso, peso, altezza),
profilo clinico (diagnosi attive, comorbidità, terapie in corso) e profilo
cardiovascolare (componenti dello score CHA$_2$DS$_2$-VASc, funzionalità
renale ed epatica~\cite{lip2010cha2ds2vasc}).
Tutti i campi sono opzionali per costruzione: l'assenza di un valore è
rappresentata come \texttt{None} Python e si propaga come \texttt{UNKNOWN}
nella logica di Kleene a tre valori (\Cref{sec:rule-based}).
La scelta di rendere tutti i campi opzionali, anziché definire un sottoinsieme
obbligatorio, risponde a un'esigenza del dominio: nella pratica clinica reale
il profilo paziente è spesso incompleto al momento della prescrizione, e
un sistema che rigettasse le richieste incomplete risulterebbe inutilizzabile
proprio nei casi più frequenti e più bisognosi di supporto.
Il campo \texttt{missing\_fields} di \texttt{EvaluationResult} costituisce
il canale strutturato attraverso cui il sistema comunica al clinico
esattamente quali informazioni aggiuntive permetterebbero di risolvere
un esito \decision{NON\_DETERMINABILE}.

Le regole normative sono codificate in un Domain-Specific Language (DSL)
in formato YAML~\cite{yamlspec}, che trasforma ogni regola in un dato
strutturato: un file di testo versionabile, leggibile da personale clinico
non tecnico e aggiornabile al ritmo degli aggiornamenti normativi senza
che l'interprete debba essere ricompilato.
Ogni regola include obbligatoriamente cinque campi: un identificatore univoco
(\texttt{rule\_id}), il tipo di regola nell'insieme
\{\texttt{SCOPE}, \texttt{PATH}, \texttt{DOSE}, \texttt{EXCEPTION},
\texttt{EXCLUSION}\}, un campo \texttt{evaluation\_order} che determina
la posizione nella pipeline, un'ancora normativa con file PDF, pagina
e estratto verbatim, e un albero di condizioni espresso tramite operatori
logici componibili (\texttt{AND}, \texttt{OR}, \texttt{NOT}, \texttt{IS\_TRUE},
\texttt{SCORE\_RANGE\_GTE} e altri).
Si è preferito YAML rispetto a JSON per la leggibilità della sintassi e
rispetto a un linguaggio di regole proprietario per la disponibilità di
tooling standard; il costo è l'assenza di tipizzazione statica,
compensata dalla validazione Pydantic all'avvio del servizio.

\subsection{Tre famiglie di regole e logica Kleene}
\label{subsec:famiglie-regole}

Le regole del DSL si articolano in tre famiglie logiche distinte, corrispondenti
ai nodi decisionali delle Note AIFA.
Le \emph{regole di inclusione} (tipo \texttt{PATH} e \texttt{INCLUSION})
verificano che il paziente soddisfi almeno uno dei criteri positivi previsti
dalla Nota: se nessun criterio di inclusione risulta soddisfatto,
l'esito è \decision{NON\_RIMBORSABILE} con short-circuit immediato.
Le \emph{regole di esclusione} (tipo \texttt{EXCLUSION}) verificano
l'assenza di controindicazioni assolute: se almeno una esclusione vale
\texttt{TRUE}, l'esito è \decision{NON\_RIMBORSABILE} con priorità
sulle inclusioni.
Le \emph{regole di eccezione} (tipo \texttt{EXCEPTION}) gestiscono
i casi di routing verso un'altra Nota o di bypass anticipato della pipeline:
il loro esito non è né rimborsabilità né non-rimborsabilità, ma
\decision{ROUTED}, che indica al clinico di rivolgersi al percorso
prescrittivo di una Nota diversa.

\begin{definition}[Processo decisionale del Rule Engine]
\label{def:decisione}
Sia $\mathcal{R}$ l'insieme ordinato delle regole applicabili a una coppia
$(\text{nota\_id}, \text{drug\_id})$, e sia
$v_\mathcal{K} \colon \text{Regola} \times \text{PatientProfile}
\to \{\texttt{TRUE}, \texttt{FALSE}, \texttt{UNKNOWN}\}$
la funzione di valutazione Kleene.
Il processo decisionale del Rule Engine è la funzione
$\delta \colon \mathcal{R} \times \text{PatientProfile} \to \mathcal{D}$,
con $\mathcal{D}$ l'insieme dei quattro esiti introdotto nella
\Cref{sec:requisiti},
definita dalla pipeline a undici fasi (numerate 0--10): la fase~0 valida
l'input e computa le variabili derivate; le fasi~1--5 corrispondono ai nodi
SCOPE, EXCEPTION, EXCL\_HARD, INCLUSION e PATHWAY, ciascuna con uscita
di short-circuit immediata; le fasi~6--8 accumulano flag clinici;
la fase~9 risolve gli eventuali conflitti fra flag di guidance;
la fase~10 applica l'invariante di sicurezza e produce l'esito finale.
Il risultato è \decision{non\_determinabile} se e solo se nessuna
regola ha prodotto un short-circuit e almeno una condizione ha
propagato \texttt{UNKNOWN}.
\end{definition}

La logica di Kleene a tre valori~\cite{kleene1952introduction} permea l'intero
processo di valutazione: ogni operatore del DSL è implementato in aritmetica
Kleene, garantendo che l'incertezza si propaghi correttamente attraverso
le composizioni logiche.
In particolare, l'operatore \texttt{SCORE\_RANGE\_GTE} opera in aritmetica
Kleene sui punteggi parziali: se uno o più componenti dello score sono
\texttt{UNKNOWN}, il confronto con la soglia non degrada a \texttt{FALSE}
ma rimane \texttt{UNKNOWN}, cosicché il sistema segnali l'indeterminazione
anziché negare erroneamente la rimborsabilità.

% ─────────────────────────────────────────────────────────────
\section{Il layer esplicativo: pipeline RAG}
\label{sec:layer-esplicativo}

La componente neurale del sistema è una pipeline RAG progettata per fornire
all'LLM una finestra di contesto strettamente ancorata al testo PDF delle
Note AIFA~\cite{lewis2020rag}.
La progettazione affronta quattro problemi specifici del dominio normativo
italiano: la natura tabellare e a clausole annidate dei PDF; il vocabolario
ibrido (terminologia latina farmacologica, acronimi inglesi, perifrasi
cliniche italiane); la necessità di recuperare chunk che mantengano l'unità
sintattica delle clausole booleane composte; e il rischio di allucinazione
nelle citazioni normative.

\subsection{Ingestione e chunking}
\label{subsec:ingestione}

La fase di ingestione, denominata \texttt{ingest\_v2}, è una fase offline
che si esegue una tantum su ogni aggiornamento dei PDF.
L'estrazione del testo grezzo è realizzata tramite PyMuPDF~\cite{pymupdf},
che preserva la struttura a pagina dei documenti e consente di associare
ogni chunk al numero di pagina sorgente, informazione indispensabile per
le ancore normative.
Il chunker adottato è \emph{sentence-aware}: il target è di $1800$ caratteri
per chunk, con un overlap di due frasi fra chunk consecutivi.
La scelta del target a $1800$ caratteri è motivata dalla finestra di $512$
token del modello di embedding: a una densità media di circa $4$ caratteri
per token sull'italiano normativo, $1800$ caratteri sono codificabili senza
troncamento.
L'overlap a livello di frase, anziché a livello di caratteri, preserva
l'unità sintattica delle clausole normative, in particolare le condizioni
booleane composte (``in pazienti con\ldots oppure\ldots''), che un overlap
fisso a caratteri spezzerebbe a metà.
Un manifest di ingestione registra il checksum SHA-256 di ogni PDF e di
ogni chunk indicizzato, garantendo la tracciabilità della base documentale.

\subsection{Modello di embedding, indice e retriever}
\label{subsec:embedding-retriever}

Il modello di embedding è \embedmodel, a $768$ dimensioni.
La scelta del multilinguismo è motivata dal vocabolario ibrido dei PDF AIFA,
che alternano terminologia latina farmacologica, acronimi clinici
(CHA$_2$DS$_2$-VASc, NAO per nuovi anticoagulanti orali, VFG per velocità di filtrazione glomerulare) e perifrasi cliniche in italiano.
Un modello monolingue italiano non gestirebbe la sottocomponente latina;
un modello generalista inglese degraderebbe la qualità sulle perifrasi
normative italiane~\cite{reimers2019sbert}.
I vettori sono indicizzati in ChromaDB con metrica coseno, organizzati
in quattro collezioni distinte (una per Nota), così da concentrare il
retrieval sulla Nota pertinente alla richiesta in esame.

Il retriever applica una strategia in due passi: recupero top-$k$ per
similarità coseno, seguito da un reranker cross-encoder \reranker{}
che riordina i candidati per rilevanza contestuale~\cite{nogueira2019passage}.
Il reranker è stato scelto specificamente per il dominio italiano, in quanto
è addestrato su coppie di frasi in italiano, e consente di distinguere
chunk tematicamente vicini ma normalmente irrilevanti (ad esempio,
tabelle di dosaggio) da chunk che contengono la clausola condizionale
pertinente al caso in esame.

\subsection{Generazione e vincolo di fedeltà}
\label{subsec:generazione}

Il modello generativo è \llmmodel, eseguito localmente tramite
Ollama~\cite{ollama2024} con parametri fissi (\texttt{temperature=0},
\texttt{seed=42}) per garantire la riproducibilità della generazione.
Il prompt è strutturato in sezioni rigide: la \emph{decisione del Rule Engine}
viene iniettata prima del contesto normativo, in una sezione marcata
\texttt{[DECISIONE DETERMINISTICA --- NON CONTRADDIRE]};
il \emph{contesto normativo} contiene i chunk recuperati, ognuno preceduto
da un'etichetta \texttt{[FONTE: \textit{pdf}, p.\,\textit{n}]};
una sezione \emph{FONTI} in coda al prompt è ricomposta deterministicamente
dal sistema, non dall'LLM.

\begin{definition}[Insieme dei chunk rilevanti]
\label{def:chunk-rilevanti}
Sia $\mathcal{C}$ l'insieme di tutti i chunk indicizzati per una Nota.
Dato un profilo paziente $p$ e un \texttt{EvaluationResult} $e$,
l'insieme dei chunk rilevanti $\mathcal{C}^*(p, e) \subseteq \mathcal{C}$
è l'output del retriever a due stadi: la prima fase seleziona i $k_1$
chunk con similarità coseno massima rispetto alla query $q(p, e)$;
la seconda fase applica il reranker \reranker{} e seleziona i $k_2 \leq k_1$
chunk con score cross-encoder più alto, dove $k_2$ è un parametro
del sistema fissato a progetto.
Il testo della spiegazione finale deve citare almeno un chunk in
$\mathcal{C}^*(p, e)$ e non può contraddire la decisione contenuta in $e$.
\end{definition}

Un controllo post-generazione, l'\emph{invariante di sicurezza},
confronta la decisione testuale prodotta dall'LLM con la decisione
deterministica del Rule Engine: in caso di divergenza, la spiegazione
viene scartata e sostituita da un template conservativo che riporta
la decisione corretta e invita il clinico a consultare direttamente
il testo normativo.
La \Cref{fig:pipeline-rag} illustra il flusso complessivo dalla richiesta
al testo di risposta.

\begin{figure}[ht]
\centering
\begin{tikzpicture}[
  phase/.style={
    rectangle, draw=black!70, thick,
    minimum width=4.4cm, minimum height=0.85cm,
    text centered, font=\small
  },
  decision/.style={
    diamond, draw=black!70, thick, aspect=2.2,
    minimum width=4.4cm, minimum height=0.75cm,
    text centered, font=\small
  },
  store/.style={
    rectangle, draw=black!50, dashed,
    minimum width=3.0cm, minimum height=0.75cm,
    text centered, font=\small
  },
  arrow/.style={-Stealth, thick},
  lbl/.style={font=\scriptsize, fill=white, inner sep=1pt},
  node distance=0.65cm and 3.2cm
]

\node[phase] (query)  {Query: $(p,\, \text{nota\_id})$};
\node[phase, below=of query] (re) {Rule Engine \texttt{/evaluate}};
\node[phase, below=of re] (embed) {Embedding query ($\embedmodel$)};
\node[phase, below=of embed] (cosine) {Retrieval coseno top-$k_1$ su ChromaDB};
\node[phase, below=of cosine] (rerank) {Reranking cross-encoder ($k_2 \leq k_1$)};
\node[phase, below=of rerank] (prompt) {Composizione prompt strutturato};
\node[phase, below=of prompt] (llm) {Generazione testo (\llmmodel, $T{=}0$)};
\node[decision, below=of llm] (check) {Invariante di sicurezza rispettato?};
\node[phase, below left=0.8cm and 1.2cm of check] (ok) {\texttt{CDSSResponse} finale};
\node[phase, below right=0.8cm and 1.2cm of check] (fallback) {Template conservativo};

\draw[arrow] (query.south)  -- (re.north);
\draw[arrow] (re.south)     -- (embed.north);
\draw[arrow] (embed.south)  -- (cosine.north);
\draw[arrow] (cosine.south) -- (rerank.north);
\draw[arrow] (rerank.south) -- (prompt.north);
\draw[arrow] (prompt.south) -- (llm.north);
\draw[arrow] (llm.south)    -- (check.north);
\draw[arrow] (check.south west) -- node[lbl,left]{Sì} (ok.north);
\draw[arrow] (check.south east) -- node[lbl,right]{No} (fallback.north);

\end{tikzpicture}
\caption{Flusso della pipeline RAG dalla richiesta al testo di risposta.
  Il Rule Engine viene invocato per primo e la sua decisione è iniettata
  nel prompt prima del contesto documentale. Il controllo dell'invariante
  di sicurezza in uscita dalla generazione garantisce che la spiegazione
  testuale non contraddica mai la decisione deterministica.}
\label{fig:pipeline-rag}
\end{figure}

% ─────────────────────────────────────────────────────────────
\section{Contratto API e flusso end-to-end}
\label{sec:orchestrazione}

L'Orchestratore è il componente di integrazione del sistema.
Esso espone verso l'esterno due endpoint: \texttt{/explain} e \texttt{/health}.
L'endpoint \texttt{/explain} riceve una \texttt{CDSSRequest} in formato JSON,
invoca il Rule Engine tramite \texttt{/evaluate} e, condizionatamente
all'esito simbolico, attiva la pipeline RAG per la generazione della
spiegazione.
L'endpoint \texttt{/health} consente al sistema di monitoraggio e al
processo di ingestione di verificare la raggiungibilità del servizio
prima di avviare operazioni critiche.

\begin{sloppypar}
Il contratto API è versionato tramite i campi \texttt{schema\_version}
(\schemaversion) e \texttt{engine\_version} (\engineversion), esposti
in ogni risposta.
La relazione fra i due campi è deliberata: \texttt{schema\_version}
identifica il contratto JSON degli oggetti scambiati, mentre
\texttt{engine\_version} identifica la versione del codice che interpreta
le regole YAML; un aggiornamento normativo che non modifica il contratto
JSON incrementa soltanto \texttt{engine\_version}, mentre un cambiamento
nello schema degli oggetti Pydantic richiede l'aggiornamento di entrambi.
Questo disaccoppiamento consente ai client (la demo, gli script di valutazione)
di gestire la compatibilità in modo esplicito.
\end{sloppypar}

La gestione dei casi \decision{NON\_DETERMINABILE} merita attenzione separata.
Quando il Rule Engine restituisce questo esito, l'Orchestratore non sopprime
la richiesta alla pipeline RAG: genera comunque una spiegazione, ma
etichettata esplicitamente come \emph{informativa}, distinta dalla spiegazione
di una decisione definitiva.
Il testo informa il clinico su quali campi mancanti impediscono la
determinazione e quali valori ipotetici, se confermati, porterebbero
a un esito positivo o negativo; questa modalità è la base della funzionalità
what-if, il quinto dei requisiti funzionali sopra elencati.

I casi \decision{ROUTED} vengono invece gestiti con un pattern di rinvio:
l'Orchestratore restituisce nella risposta l'identificatore della Nota
verso cui il caso è stato instradato, lasciando al client il compito
di reindirizzare la richiesta al percorso prescrittivo corretto.
Il testo esplicativo accompagna il routing con le motivazioni normative
(ad esempio: il farmaco prescritto non rientra nell'ambito della \notezero\
ma potrebbe rientrare nella \notetredici\ per indicazioni ipolipemizzanti).

Sul piano della robustezza, l'Orchestratore implementa un meccanismo di
timeout esplicito per le chiamate a Ollama: se il modello non risponde
entro il limite configurato, viene restituito il template conservativo
dell'invariante di sicurezza, che almeno garantisce al clinico la decisione
simbolica corretta anche in assenza della spiegazione testuale.

L'architettura sviluppata in questo capitolo si regge su un solo principio:
la separazione netta fra il layer decisionale simbolico, formalmente
verificabile e deterministico, e il layer esplicativo neurale, utile e ricco
ma intrinsecamente stocastico. Il rule engine garantisce la correttezza
normativa attraverso una pipeline a fasi ordinate con logica Kleene, regole
YAML versionabili e ancore normative verificabili; la pipeline RAG garantisce
la qualità della spiegazione attraverso un chunker sentence-aware, un
retriever a due stadi e un controllo post-generazione che protegge il
clinico da eventuali allucinazioni del modello. Il contratto fra i due
layer è espresso da tipi Pydantic validati a runtime, e il sistema è
integralmente deployato su hardware locale per rispettare i vincoli di
privacy dei dati clinici. Questa separazione è anche ciò che rende
nette le domande di ricerca: poiché la decisione non dipende mai dal modello
generativo, RQ1 può chiedere se il solo layer simbolico decida come l'annotatore
umano, e RQ3 e RQ6 possono interrogare la fedeltà e la verificabilità della
spiegazione senza che un loro eventuale fallimento comprometta la correttezza
dell'esito. Resta da mostrare come queste decisioni di
progetto si traducono in codice eseguibile e riproducibile, dalla funzione
\texttt{evaluate} ai parametri di generazione del modello linguistico: è
il compito del \Cref{ch:implementazione}.
